\label{ch:geometrie_vs_info}
% Source document: 243

\epigraph{That FFGFT can be mapped into an information formalism
does not prove that information is the prior. It shows that both
descriptions are projections of the same structure.}{Doc.\,243}

\section{The Fundamental Question}

Is information or geometry the more fundamental primitive? This question is
controversially discussed in the IPI context --- Onur Teker (UIFT) and others advocate
information as the ontological foundation from which geometry \emph{follows}.

FFGFT gives a precise answer: Both descriptions are
\textbf{projections of the same structure}. Which one chooses does not decide
the physics --- but it does decide what is lost in the projection.

\section{The Geometric State Space in the FFGFT}

The fundamental state space in FFGFT is the $T^4$ torus with
winding numbers $(n_\theta, n_\varphi) \in \mathbb{Z}^2$ and phase
$\theta_4 \in [0, 2\pi)$. A physical state $\Psi$ is
completely characterized by:
\[
  \Psi \;=\; \bigl(n_\theta,\; n_\varphi,\; \theta_4,\; \varepsilon\bigr)
\]
where $\varepsilon \approx 0{,}001$ is the non-closure (Doc.\,193).

\section{The Information Projection \texorpdfstring{$\pi_I$}{pi\_I}}

The projection $\pi_I$ maps the geometric state onto an
information representation:
\[
  \pi_I:\; \Psi(n_\theta, n_\varphi, \theta_4, \varepsilon)
  \;\longmapsto\;
  \bigl(I(n_\theta, n_\varphi),\;\operatorname{sgn}(\theta_4),\;
        \langle\varepsilon\rangle_\text{therm}\bigr)
\]
where $I = \log_2|\mathcal{W}|$ is the number of bits of admissible
winding configurations.

\textbf{Decisive:} $|\mathcal{W}|$ must be finite for $I$
to be well-defined. The $T^4$ geometry provides this bound via the maximum
admissible winding number $n_\text{max} \sim 1/\xi$. Without this geometric
bound, $I$ diverges --- the bit counting \emph{presupposes} the geometric
state space, it does not create it.

\section{What is Lost in the Projection}

The mapping $\pi_I$ is not a bijection. Three pieces of information are lost:

\begin{enumerate}
  \item \textbf{Continuous phase} $\theta_4 \in [0, 2\pi)$ is reduced to binary
    $\operatorname{sgn}(\theta_4)$ --- loss of a real degree of freedom.
  \item \textbf{Winding numbers} $(n_\theta, n_\varphi)$ are projected onto a scalar bit number
    $I$ --- loss of directional information.
  \item \textbf{Non-closure} $\varepsilon$ is thermally averaged ---
    loss of deterministic fine structure.
\end{enumerate}

\begin{keyresult}
Information \emph{as such} is not an ontological primitive in FFGFT ---
it is a useful description that emerges from the
geometric state space during projection. Wheeler's "It from Bit" reverses the
derivation: FFGFT says "Bit from It" --- where "It" is the $T^4$ torus.
\end{keyresult}

\section{Consequence for the Landauer Discussion}

That $\pi_I$ is not a bijection has a direct consequence for
Landauer: The erasure of a bit corresponds to a region reduction in the
geometric state space, not in an abstract information space.
Thermodynamics counts winding states --- it does not know bits.
Bits are a descriptive decision about this state space.

\textit{Open Question} (Doc.\,243): Is there a complete reconstruction
$\pi_I^{-1}$ that uniquely infers the geometric state from the
information representation? If not --- and this is the
FFGFT thesis --- then geometry is the deeper primitive.